\begin{ouverture}[geometrie]{14}{Médiatrice et équidistance}{Partie 3 · Espace et géométrie}
{Perpendiculaire · Milieu · Points équidistants}
\objectif{Connaître la définition de la médiatrice d’un segment et la construire.}
\objectif{Utiliser la propriété caractéristique dans les deux sens.}
\objectif{Résoudre des problèmes de points à égale distance.}
\prerequis{Milieu d’un segment ; perpendiculaires ; compas.}
\end{ouverture}
\pagestyle{otMain}\essentiel\fontsize{11}{15}\selectfont
\phantomsection\label{cours-14-1}
\essSection{1}{Définir une médiatrice}
\begin{definition}La médiatrice d’un segment est la droite perpendiculaire à ce segment qui passe par son milieu.\end{definition}
\label{fig-cours-14-1-cours-mediatrice}\figureGeo{cours-mediatrice}

\essSection{2}{Construire avec les instruments}
\begin{savoirfaire}[À la règle et à l’équerre]On place le milieu I du segment [AB], puis on trace la droite perpendiculaire au segment passant par I.\end{savoirfaire}
\begin{savoirfaire}[Au compas]Avec un même écartement, plus grand que la moitié de AB, on trace un arc de centre A et un arc de centre B, au-dessus puis au-dessous du segment. Les deux points d’intersection des arcs sont sur la médiatrice : on trace la droite qui les joint.\end{savoirfaire}
\clearpage
\phantomsection\label{cours-14-2}
\essSection{3}{Deux sens d’une caractérisation}
\begin{propriete}Si un point appartient à la médiatrice d’un segment [AB], alors il est à égale distance de A et de B. Réciproquement, si un point est à égale distance de A et de B, alors il appartient à la médiatrice de [AB].\end{propriete}
\label{fig-cours-14-2-cours-equidistance}\figureGeo{cours-equidistance}

\begin{exemple}Sur la figure, M est sur la médiatrice de [AB] et $MA=4$ cm : on peut affirmer, sans mesurer, que $MB=4$ cm. En revanche $NA\ne NB$ : N n’est pas sur la médiatrice.\end{exemple}
\essSection{4}{Raisonner avec des lieux de points}
\begin{savoirfaire}[Rédiger en trois temps]Je sais que $PA=PB=5$ cm (donnée). Or un point à égale distance de A et de B est sur la médiatrice de [AB] (propriété). Donc P est sur la médiatrice de [AB] (conclusion).\end{savoirfaire}
\begin{attention}Une impression visuelle ne suffit pas : des longueurs qui « semblent » égales doivent être données, mesurées ou justifiées.\end{attention}

\clearpage\rubrique{Découverte}
\begin{blocLivre}{ex-14-1}\begin{exercice}[1]{Reconnaître la médiatrice}
Le segment [AB] mesure 6 cm. Parmi les droites $(d_1)$, $(d_2)$, $(d_3)$ et $(d_4)$, laquelle est la médiatrice de [AB] ? Justifier en éliminant les autres.
\label{fig-ex-14-1}\figureGeo{med-1}
\end{exercice}\end{blocLivre}
\begin{blocLivre}{ex-14-2}\begin{exercice}[1]{Coder une médiatrice}
La droite (d) est la médiatrice du segment [AB]. Reproduire la figure et ajouter tous les codages utiles.
\label{fig-ex-14-2}\figureGeo{med-2}
\end{exercice}\end{blocLivre}
\begin{blocLivre}{ex-14-3}\begin{exercice}[1]{Construire à l’équerre}
Tracer un segment [AB] de 6 cm, puis sa médiatrice à la règle graduée et à l’équerre.
\label{fig-ex-14-3}\figureGeo{med-3}
\end{exercice}\end{blocLivre}
\begin{blocLivre}{ex-14-4}\begin{exercice}[1]{Une longueur sans mesurer}
Le point M est sur la médiatrice du segment [AB] et $MA=4{,}2$ cm. Que vaut MB ? Justifier.
\label{fig-ex-14-4}\figureGeo{med-4}
\end{exercice}\end{blocLivre}
\begin{blocLivre}{ex-14-5}\begin{exercice}[1]{Sur la médiatrice ?}
On sait que $PA=PB=3{,}5$ cm, que $QA=3$ cm et que $QB=4$ cm. Lesquels des points P et Q sont sur la médiatrice de [AB] ?
\label{fig-ex-14-5}\figureGeo{med-5}
\end{exercice}\end{blocLivre}

\rubrique{Entraînement}
\begin{blocLivre}{ex-14-6}\begin{exercice}[2]{Construire au compas}
Tracer un segment [CD] de 5 cm. Construire sa médiatrice au compas, avec un écartement de 4 cm.
\label{fig-ex-14-6}\figureGeo{med-6}
\end{exercice}\end{blocLivre}
\begin{blocLivre}{ex-14-7}\begin{exercice}[2]{Le triangle isocèle}
Le triangle ABC vérifie $AB=AC=5$ cm et $BC=6$ cm. Démontrer que A est sur la médiatrice de [BC]. Où cette médiatrice coupe-t-elle le segment [BC] ?
\label{fig-ex-14-7}\figureGeo{med-7}
\end{exercice}\end{blocLivre}
\begin{blocLivre}{ex-14-8}\begin{exercice}[2]{Le losange}
ABCD est un losange : ses quatre côtés ont la même longueur. Démontrer que la droite (BD) est la médiatrice du segment [AC].
\label{fig-ex-14-8}\figureGeo{med-8}
\end{exercice}\end{blocLivre}
\begin{blocLivre}{ex-14-9}\begin{exercice}[2]{Un point sur une droite}
Construire le point M de la droite (d) qui est à égale distance de A et de B.
\label{fig-ex-14-9}\figureGeo{med-9}
\end{exercice}\end{blocLivre}
\begin{blocLivre}{ex-14-10}\begin{exercice}[2]{Des cercles passant par A et B}
Les points A et B sont distants de 4 cm. Où sont situés les centres des cercles qui passent à la fois par A et par B ? En tracer deux.
\label{fig-ex-14-10}\figureGeo{med-10}
\end{exercice}\end{blocLivre}

\rubrique{Problèmes}
\begin{blocLivre}{ex-14-11}\begin{exercice}[3]{L’arrêt de bus}
Les maisons A et B sont situées de part et d’autre d’une route droite (d). On veut placer un arrêt de bus S sur la route, à la même distance des deux maisons. Construire S.
\label{fig-ex-14-11}\figureGeo{med-11}
\end{exercice}\end{blocLivre}
\begin{blocLivre}{ex-14-12}\begin{exercice}[3]{Plus près de A ou de B ?}
Colorier la zone des points plus proches de A que de B. Quelle ligne sépare les deux zones ?
\label{fig-ex-14-12}\figureGeo{med-12}
\end{exercice}\end{blocLivre}
\begin{blocLivre}{ex-14-13}\begin{exercice}[3]{Le pliage}
On plie une feuille de façon que le point A vienne exactement sur le point B. Expliquer pourquoi le pli est la médiatrice du segment [AB].
\label{fig-ex-14-13}\figureGeo{med-13}
\end{exercice}\end{blocLivre}
\begin{blocLivre}{ex-14-14}\begin{exercice}[3]{Raisonner}
La droite (d) est la médiatrice de [AB]. Le point C est sur (d) et $CA=5$ cm. Le point D vérifie $DA=DB=3{,}5$ cm. Que vaut CB ? Le point D est-il sur (d) ? Rédiger les deux réponses.
\label{fig-ex-14-14}\figureGeo{med-14}
\end{exercice}\end{blocLivre}
\begin{blocLivre}{ex-14-15}\begin{exercice}[3]{Vrai ou faux}
Léo affirme : « Un point à égale distance de A et de B est toujours le milieu de [AB]. » A-t-il raison ?
\label{fig-ex-14-15}\figureGeo{med-15}
\end{exercice}\end{blocLivre}

\clearpage
